\documentclass[11pt,a4paper]{article}
\usepackage[hyperref]{acl}
\usepackage{times}
\usepackage{latexsym}
\usepackage{graphicx}
\usepackage{booktabs}
\usepackage{amsmath}
\usepackage{multirow}
\usepackage{xcolor}
\usepackage{url}

\aclfinalcopy

\title{Not All Policy Ambiguity Is Equal: Differential Vulnerability Profiles\\ Across LLM Agent Families}

\author{Anonymous}

\date{}

\begin{document}
\maketitle

% ============================================================
% ABSTRACT
% ============================================================
\begin{abstract}
Large Language Model (LLM) agents increasingly operate under natural-language policy specifications that inevitably contain ambiguity.
Yet no prior work causally quantifies how policy ambiguity affects agent safety, nor whether different ambiguity types pose different risks.
We introduce a matched-pair causal design over a six-type, two-layer ambiguity taxonomy and evaluate 2{,}997 agent episodes across five architecturally diverse models (GPT-5.4, GPT-4.1, Claude~Sonnet~4.6, Qwen3-235B, DeepSeek-v3) on the $\tau^2$-bench platform.
Ambiguous policy clauses increase safety-violation rates by 34.0 percentage points (pp) over unambiguous matched variants (OR\,=\,15.1, $p<$0.0001), an effect that replicates across all five models.
A perturbation control confirms that non-ambiguity surface changes produce near-zero effect ($\Delta$\,=\,$-$0.67pp), isolating ambiguity-specific causation.
Within the ambiguous condition, a three-tier type hierarchy emerges: \emph{incompleteness} is the most dangerous (53.6\%), while \emph{coreferential} ambiguity is the least (31.6\%).
Critically, models exhibit dramatically different vulnerability profiles---authorization-scope violation rates range from 12\% (Claude) to 64\% (GPT-5.4), revealing that safety training differentially mitigates specific ambiguity types.
Partial cross-benchmark replication on SOPBench (healthcare and banking domains) confirms the binary effect ($\Delta$\,=\,18.3pp, $p<$0.0001) while revealing that effect size and type severity ordering are moderated by domain and benchmark strictness.
We release all 300 clause pairs, evaluation scripts, and analysis code.
\end{abstract}

% ============================================================
% 1. INTRODUCTION
% ============================================================
\section{Introduction}

When a customer returns a customized item within the return window, should the agent grant a full refund?
A policy stating ``items may be returned within 30 days'' leaves the scope of ``items'' ambiguous---does it include custom orders?
In our experiments, agents operating under this ambiguous clause grant unauthorized refunds, while the same agents correctly refuse when the policy specifies ``standard items only.''
One word of ambiguity triggers a \$200 safety violation.

LLM-based agents are increasingly deployed to execute tasks governed by natural-language policy specifications: customer-service protocols, healthcare guidelines, and banking procedures.
These specifications, written for human operators, inevitably contain linguistic and structural ambiguity \cite{berry2003ambiguity}.
While ambiguity in requirements engineering has been studied extensively \cite{ferrari2019nlp}, its \emph{safety} impact on autonomous agents remains unquantified.
Prior benchmarks evaluate agent compliance \cite{li2025sopbench}, robustness to adversarial inputs, or instruction-following fidelity, but none isolate the causal effect of policy ambiguity on safety violations at the granularity of ambiguity \emph{types}.

Berry and Kamsties~(\citeyear{berry2003ambiguity}) distinguish \emph{nocuous} ambiguity (divergent interpretations causing harm) from \emph{innocuous} ambiguity (convergent interpretations despite ambiguity).
In human requirements engineering, shared organizational context renders most ambiguity innocuous.
LLM agents lack this shared context.
Are policy ambiguities systematically nocuous for agents?
Does the \emph{type} of ambiguity predict the severity of harm?
Does the answer depend on the \emph{model}?

We present the first causal, type-granularity, cross-model study of policy ambiguity and agent safety.
Our contributions are:
\begin{enumerate}
\item \textbf{Causal quantification}: A matched-pair design with perturbation control showing that ambiguous policy clauses increase safety-violation rates by +34.0pp (OR\,=\,15.1) across five models, while non-ambiguity perturbations produce near-zero effect ($\Delta$\,=\,$-$0.67pp).
\item \textbf{Differential Nocuity}: A three-tier type hierarchy (incompleteness $\gg$ mid cluster $\gg$ coreferential) and a model$\times$type interaction where authorization-scope vulnerability varies 52pp across models, revealing that safety training differentially addresses specific ambiguity types while incompleteness remains a universal blind spot.
\item \textbf{Ambiguity Blindness}: Error analysis of 614 violation traces showing $<$5\% meta-cognitive awareness of ambiguity, $>$88\% surface adoption, and $<$1\% user escalation, with type-specific failure signatures.
\end{enumerate}

% ============================================================
% 2. RELATED WORK
% ============================================================
\section{Related Work}

\paragraph{Ambiguity in NLP and Requirements Engineering.}
Ambiguity detection has a long history in requirements engineering (RE).
Chantree et al.~(\citeyear{chantree2006identifying}) developed the first computational approach to nocuous ambiguity detection.
Ferrari et al.~(\citeyear{ferrari2019nlp}) extended NLP-based cross-domain ambiguity detection.
More recently, Li et al.~(\citeyear{li2024drip}) introduced DRIP-R, classifying five RE ambiguity types and showing that LLMs diverge in their interpretations---but did not measure whether this divergence leads to safety violations.
Massey et al.~(\citeyear{massey2014evaluating}) proposed a six-type taxonomy for policy ambiguities.
Our work builds on these taxonomies but shifts from \emph{interpretation divergence} to \emph{safety-violation causation}.

\paragraph{Agent Safety Evaluation.}
Recent benchmarks evaluate agent safety from different angles: PhantomPolicy tests compliance with hidden policies, PolicyBank examines policy gap types, and AgentPex probes procedural errors.
$\tau^2$-bench \cite{tau2bench2025} provides a multi-turn agent evaluation platform with grounded safety metrics.
SOPBench \cite{li2025sopbench} evaluates SOP-following across healthcare and banking domains with rule-based verification.
None of these benchmarks isolate ambiguity type as a causal variable or provide matched-pair designs.

\paragraph{LLM Robustness and Prompt Sensitivity.}
Sclar et al.~(\citeyear{sclar2024quantifying}) demonstrate that minor prompt variations significantly affect LLM performance.
Our perturbation control (\S\ref{sec:perturbation}) distinguishes ambiguity-specific effects from general prompt brittleness, showing that non-ambiguity surface changes produce negligible effect.

% ============================================================
% 3. AMBIGUITY TAXONOMY
% ============================================================
\section{Ambiguity Taxonomy}
\label{sec:taxonomy}

We define a two-layer, six-type taxonomy grounded in Li et al.~(\citeyear{li2024drip}) and Massey et al.~(\citeyear{massey2014evaluating}).

\paragraph{Specification Layer.}
Ambiguities arising from gaps or conflicts in the policy's \emph{content}:
\begin{itemize}
\item \textbf{Incompleteness}: Missing conditions, exceptions, or boundary specifications. Disambiguation requires adding information not present in the original text.
\item \textbf{Authorization Scope}: Ambiguous permission or role boundaries that are stated but vague (e.g., ``authorized personnel'' without specifying functions).
\item \textbf{Conditional Precedence}: Conflicting or underspecified rule ordering when multiple conditions apply simultaneously.
\end{itemize}

\paragraph{Linguistic Layer.}
Ambiguities arising from the \emph{language} used to express the policy:
\begin{itemize}
\item \textbf{Lexical}: Word-level polysemy or homonymy (e.g., ``interest'' in banking contexts).
\item \textbf{Scopal}: Quantifier, negation, or modal scope interactions (e.g., ``all items are not eligible'' --- universal negation or negated universal?).
\item \textbf{Coreferential}: Ambiguous pronoun or relative-clause attachment creating referential confusion.
\end{itemize}

The two layers differ in their \emph{disambiguation channels}: linguistic ambiguity can sometimes be partially resolved via tool responses, world knowledge, or conversational context; specification ambiguity has no external resolution channel, forcing agents to guess.

% ============================================================
% 4. EXPERIMENTAL DESIGN
% ============================================================
\section{Experimental Design}
\label{sec:method}

\subsection{Clause Generation}
\label{sec:clause_gen}

We construct 300 matched clause pairs (50 per type $\times$ 6 types), each consisting of an \emph{ambiguous} variant and an \emph{unambiguous} variant that conveys the same policy intent.
Table~\ref{tab:clause_gen} summarizes the generation statistics.
Specification-layer clauses require significantly more tokens to disambiguate than linguistic-layer clauses (mean 34.0 vs.\ 14.8 tokens; $t$-test $p<$0.001), reflecting the structural nature of specification-layer ambiguity.
Conditional-precedence clauses use a 60/40 mix of explicitly labeled and naturalistic formats.
Quality review achieved a 92\% pass rate, with 11 clauses (primarily \emph{authorization\_scope} and \emph{incompleteness}) requiring revision.

\begin{table}[t]
\centering
\small
\begin{tabular}{llccc}
\toprule
\textbf{Type} & \textbf{Layer} & \textbf{$n$} & \textbf{Avg $\Delta$ tok} & \textbf{Rev.} \\
\midrule
auth\_scope & spec & 50 & 36.1 & 3 \\
cond\_prec & spec & 50 & 31.4 & 1 \\
incompleteness & spec & 50 & 34.6 & 4 \\
lexical & ling & 50 & 18.2 & 2 \\
coreferential & ling & 50 & 7.8 & 1 \\
scopal & ling & 50 & 18.3 & 0 \\
\bottomrule
\end{tabular}
\caption{Clause pair generation summary. \textbf{Avg $\Delta$ tok}: mean token difference between ambiguous and unambiguous variants. \textbf{Rev.}: clauses requiring revision.}
\label{tab:clause_gen}
\end{table}

\subsection{Mini-Pilot Validation}

We validate two new types (\emph{coreferential}, \emph{authorization\_scope}) with a mini-pilot: 5 clauses $\times$ 2 conditions $\times$ 1 model (GPT-5.4) = 20 episodes.
Authorization scope shows a 60\% ambiguous violation rate versus 10\% unambiguous baseline; coreferential shows 20\% versus 0\%.
Two unambiguous authorization-scope clauses with residual boundary ambiguity are revised before the full study.

\subsection{Models and Platform}

We evaluate five architecturally diverse models spanning three providers: GPT-5.4 and GPT-4.1 (OpenAI, proprietary dense), Claude Sonnet 4.6 (Anthropic, proprietary dense), Qwen3-235B (Alibaba, open-source MoE), and DeepSeek-v3 (DeepSeek, open-source dense).
All experiments use the $\tau^2$-bench platform with its built-in multi-turn agent interaction and safety-violation grading.
Cross-judge evaluation uses a separate LLM judge, with each model judged by a different model to avoid self-evaluation bias.

\subsection{Full Study}

The full study comprises 300 clauses $\times$ 2 conditions $\times$ 5 models = 3{,}000 target episodes.
API models (GPT-5.4, GPT-4.1, Claude) contribute 1{,}800 episodes (1{,}793 valid, 7 failed due to API timeouts).
Open-source models (Qwen3, DeepSeek) contribute 1{,}200 episodes (1{,}188 valid, 12 failed).
Total valid episodes: 2{,}981 (99.4\% completion rate).

\subsection{Measurement Reliability}
\label{sec:reliability}

To quantify API sampling noise, we run a reliability test: 30 clauses $\times$ 3 replicates $\times$ 2 conditions $\times$ 2 models = 360 episodes.
Overall ICC = 0.919 [0.892, 0.940] (excellent).
Linguistic-layer types show higher reproducibility than specification-layer types (avg ICC 0.949 vs.\ 0.835), consistent with the disambiguation channel framework: linguistic ambiguity elicits more stable (mis)resolution patterns, while specification ambiguity forces variable guessing across runs.
Per-type ICC ranges from 0.807 (\emph{incompleteness}) to 0.989 (\emph{scopal}).
Binary ICC = 0.847 (good).

\begin{table}[t]
\centering
\small
\begin{tabular}{llcc}
\toprule
\textbf{Type} & \textbf{Layer} & \textbf{ICC} & \textbf{95\% CI} \\
\midrule
scopal & ling & 0.989 & [0.968, 0.994] \\
lexical & ling & 0.981 & [0.948, 0.987] \\
coreferential & ling & 0.976 & [0.941, 0.984] \\
auth\_scope & spec & 0.857 & [0.793, 0.896] \\
cond\_prec & spec & 0.856 & [0.783, 0.890] \\
incompleteness & spec & 0.807 & [0.728, 0.851] \\
\bottomrule
\end{tabular}
\caption{Per-type ICC (intraclass correlation) for measurement reliability. Linguistic-layer types show consistently higher reproducibility.}
\label{tab:icc}
\end{table}

\subsection{Judge Calibration}
\label{sec:judge_cal}

We calibrate the LLM judge against human ground truth using 100 stratified episodes annotated by two human experts.
Human inter-annotator agreement is $\kappa_w$\,=\,0.762 (substantial).
Human--LLM agreement is $\kappa_w$\,=\,0.807 (excellent), \emph{exceeding} human--human agreement, validating the LLM judge.
Cross-judge (LLM--LLM) agreement is $\kappa_w$\,=\,0.685 (substantial), with binary $\kappa$\,=\,0.74.
Judge--human binary F1 = 0.82.
While these results validate the measurement chain, we note that the task exhibits inherent subjectivity, particularly for authorization-scope boundary judgments.

\subsection{Realism Validation}

We compare the linguistic properties of our 300 synthetic clause pairs against 87 naturally occurring ambiguities from DRIP-R.
Three of four distributional metrics show no significant difference: Flesch--Kincaid grade (12.4 vs.\ 11.8, KS $p$\,=\,0.31), sentence length (24.7 vs.\ 22.1, $p$\,=\,0.08), and type--token ratio (0.72 vs.\ 0.69, $p$\,=\,0.15).
Parse depth is marginally higher in synthetic clauses (5.8 vs.\ 5.2, $p$\,=\,0.04), reflecting controlled subordinate-clause insertion for disambiguation.

% ============================================================
% 5. RESULTS
% ============================================================
\section{Results}
\label{sec:results}

\subsection{C1: Binary Effect}
\label{sec:c1}

Ambiguous policy clauses cause massive safety-violation increases across all five models (Table~\ref{tab:c1}).
The pooled effect is +34.0pp (OR\,=\,15.1, 95\% CI [10.78, 21.17], bootstrap CI [31.2, 36.9]pp, $p<$0.0001).
All five models individually show highly significant effects (all $p<10^{-14}$), ranging from +25.7pp (Claude) to +39.0pp (DeepSeek).
Logistic regression confirms ambiguity condition as the dominant predictor ($\beta$\,=\,2.278, $z$\,=\,19.65) after controlling for token delta, model, and type; token delta is significant but secondary ($\beta$\,=\,0.022, $z$\,=\,3.49).

\begin{table}[t]
\centering
\small
\begin{tabular}{lcccc}
\toprule
\textbf{Model} & \textbf{$\Delta$pp} & \textbf{OR} & \textbf{$p$} & \textbf{Boot. CI} \\
\midrule
Claude & 25.7 & 8.34 & $<$1e-17 & [20.9, 30.5] \\
Qwen3 & 32.4 & 12.47 & $<$1e-17 & [27.4, 37.4] \\
GPT-4.1 & 34.7 & 16.21 & $<$1e-17 & [29.7, 39.7] \\
GPT-5.4 & 38.3 & 19.53 & $<$1e-17 & [33.3, 43.3] \\
DeepSeek & 39.0 & 17.82 & $<$1e-17 & [34.0, 44.0] \\
\midrule
\textbf{Pooled} & \textbf{34.0} & \textbf{15.1} & $<$\textbf{0.0001} & \textbf{[31.2, 36.9]} \\
\bottomrule
\end{tabular}
\caption{C1: Binary effect of policy ambiguity per model. $\Delta$pp = ambiguous $-$ unambiguous violation rate. OR = McNemar odds ratio.}
\label{tab:c1}
\end{table}

Cross-benchmark replication on SOPBench (healthcare and banking domains, 982 valid episodes) confirms the binary effect ($\Delta$\,=\,18.3pp, OR\,=\,3.31, $p<$0.0001) with a smaller effect size, attributable to the stricter rule-based verifier, higher unambiguous baselines, and domain moderation (healthcare $\Delta$\,=\,16.2pp vs.\ banking $\Delta$\,=\,20.4pp).
Full SOPBench results appear in Appendix~\ref{app:sopbench}.

\subsection{Perturbation Control}
\label{sec:perturbation}

To isolate ambiguity-specific causation from general prompt sensitivity, we construct 30 non-ambiguity perturbation pairs (synonym substitution, clause reordering) and run 300 episodes.
The perturbation effect is $\Delta$\,=\,$-$0.67pp (bootstrap CI [$-$8.0, 6.67], McNemar $p$\,=\,1.0), indistinguishable from zero.
Net ambiguity-specific effect: +34.67pp.
This confirms that the +34.0pp binary effect is attributable to ambiguity, not surface-level variation.

\begin{table}[t]
\centering
\small
\begin{tabular}{lccc}
\toprule
\textbf{Metric} & \textbf{Value} & \textbf{Threshold} & \textbf{Verdict} \\
\midrule
$\Delta_{\text{perturb}}$ & $-$0.67pp & $<$5pp & PASS \\
CI upper bound & 6.67pp & $<$10pp & PASS \\
Effect ratio & $\approx$0 & $<$0.15 & PASS \\
$\Delta_{\text{net}}$ & 34.67pp & $>$20pp & PASS \\
\bottomrule
\end{tabular}
\caption{Perturbation control results. Non-ambiguity surface changes produce near-zero effect, confirming ambiguity-specific causation.}
\label{tab:perturbation}
\end{table}

\subsection{C2: Type Hierarchy}
\label{sec:c2}

Ambiguity types differ significantly in their violation rates (omnibus $\chi^2$\,=\,28.02, $p$\,=\,3.6$\times$10$^{-5}$, Cram\'er's $V$\,=\,0.137, power\,=\,0.993).
Table~\ref{tab:c2} shows the three-tier structure: \emph{incompleteness} (53.6\%) stands alone as the most dangerous type, four types cluster in a mid tier (36.8--42.8\%), and \emph{coreferential} (31.6\%) is the least dangerous.

While $V$\,=\,0.137 indicates a small effect, TOST equivalence testing confirms that types are not interchangeable.
At $\pm$15pp margins, 7/15 pairwise comparisons are significantly non-equivalent; at $\pm$10pp, 11/15; at $\pm$5pp, 14/15.
The hierarchy is real but driven by the extremes: the 22pp gap between \emph{incompleteness} and \emph{coreferential}.

Partial cross-benchmark replication on SOPBench confirms \emph{incompleteness} as the most dangerous type (39.2\%) but shows domain-sensitive mid-tier ordering ($\rho$\,=\,0.40 between $\tau^2$-bench and SOPBench rankings).
Authorization scope drops from rank 3 to rank 4 in structured healthcare and banking domains.
We therefore characterize the hierarchy as robust in customer-service contexts and partially replicated cross-domain: \emph{incompleteness} is universally dangerous, but mid-tier ordering is domain-moderated.

\begin{table}[t]
\centering
\small
\begin{tabular}{llccc}
\toprule
\textbf{Type} & \textbf{Layer} & \textbf{Viol. \%} & \textbf{Tier} & \textbf{95\% CI} \\
\midrule
incompleteness & spec & 53.6 & high & [48.4, 58.8] \\
lexical & ling & 42.8 & mid & [37.6, 48.0] \\
auth\_scope & spec & 41.6 & mid & [36.4, 46.8] \\
scopal & ling & 39.4 & mid & [34.2, 44.6] \\
cond\_prec & spec & 36.8 & mid & [31.7, 41.9] \\
coreferential & ling & 31.6 & low & [26.5, 36.7] \\
\bottomrule
\end{tabular}
\caption{C2: Type hierarchy (ambiguous condition, pooled across 5 models). $\chi^2$\,=\,28.02, $p$\,=\,3.6$\times$10$^{-5}$, $V$\,=\,0.137.}
\label{tab:c2}
\end{table}

\subsection{C4: Model $\times$ Type Interaction}
\label{sec:c4}

While the binary effect is universal, models exhibit dramatically different vulnerability profiles (model main effect: $\chi^2$\,=\,31.02, $p$\,=\,3.0$\times$10$^{-6}$, $V$\,=\,0.102).
Table~\ref{tab:c4} presents the full interaction matrix.

\begin{table}[t]
\centering
\small
\begin{tabular}{lcccccc}
\toprule
\textbf{Type} & \textbf{G5.4} & \textbf{G4.1} & \textbf{Cla.} & \textbf{Qw3} & \textbf{DS} & \textbf{Range} \\
\midrule
auth\_scope & 64 & 46 & 12 & 26 & 60 & 52pp \\
incomp. & 46 & 58 & 48 & 62 & 54 & 16pp \\
lexical & 42 & 34 & 32 & 52 & 54 & 22pp \\
coref. & 18 & 32 & 22 & 42 & 44 & 26pp \\
scopal & 42 & 38 & 32 & 37 & 48 & 16pp \\
cond\_prec & 34 & 36 & 36 & 40 & 38 & 6pp \\
\bottomrule
\end{tabular}
\caption{C4: Model $\times$ type interaction (ambiguous condition, violation \%). G5.4 = GPT-5.4, Cla. = Claude, Qw3 = Qwen3, DS = DeepSeek. Range = max $-$ min across models.}
\label{tab:c4}
\end{table}

Three key interaction patterns emerge:

\paragraph{Authorization scope is the model-differentiating type.}
The 52pp range (Claude 12\% to GPT-5.4 64\%) is the largest cross-model gap for any type.
Claude's RLHF appears to specifically address authorization-boundary awareness.
The directionality is entirely toward escalation: in authorization-scope violations, 69.2\% involve unauthorized escalation and 0\% involve over-restriction (binomial $p$\,=\,5.5$\times$10$^{-5}$), consistent across 4/5 models.

\paragraph{Incompleteness is a safety-training blind spot.}
All models show high vulnerability (46--62\%, 16pp range), indicating that safety training does not help when information is literally absent from the policy.

\paragraph{Conditional precedence is the most model-invariant type.}
The 6pp range (34--40\%) suggests that explicit rule conflicts are processed similarly by all architectures.

Inter-model agreement on \emph{which specific clauses} trigger violations is only fair (Fleiss $\kappa$\,=\,0.384): the binary effect is universal, but the specific failure triggers are model-dependent.
We characterize the interaction as domain-robust for extreme types (incompleteness universally high, conditional precedence uniformly moderate) and domain-sensitive for mid-tier types (authorization scope, coreferential).

% ============================================================
% 6. ERROR ANALYSIS
% ============================================================
\section{Error Analysis}
\label{sec:error_analysis}

We code 614 violation traces across all five models to characterize failure mechanisms (Table~\ref{tab:error}).

\begin{table}[t]
\centering
\small
\begin{tabular}{llcl}
\toprule
\textbf{Type} & \textbf{Dominant Mode} & \textbf{Aware \%} & \textbf{Mechanism} \\
\midrule
auth\_scope & unauth. escalation & 3.2 & spec $\to$ escalate \\
incomp. & assumption-based & 4.8 & spec $\to$ guess \\
scopal & scope misappl. & 5.1 & ling $\to$ misparse \\
lexical & assumption-based & 6.8 & ling $\to$ guess \\
cond\_prec & mixed & 11.2 & spec $\to$ mixed \\
coref. & flat distribution & 8.5 & ling $\to$ diverse \\
\bottomrule
\end{tabular}
\caption{Failure mode summary (614 coded violations). Dominant modes are majority-level classifications. Aware \% = proportion of traces showing meta-cognitive awareness of ambiguity.}
\label{tab:error}
\end{table}

\paragraph{Ambiguity Blindness.}
Agents show near-zero meta-cognitive awareness of policy ambiguity.
Across all types and models, the ambiguity awareness rate is 4.6\% [2.8, 6.4]; surface adoption (proceeding without noting ambiguity) is 88\% [85, 91]; and user escalation (asking for clarification) is $<$1\% [0, 1.5].
Claude shows modestly higher awareness (11\% vs.\ $<$5\% for other models), consistent with its lower violation rate.

\paragraph{Layer-Specific Failure Mechanisms.}
Specification-layer types produce ``guess-and-act'' failures: agents fabricate missing information (\emph{incompleteness}), assume broader permissions (\emph{authorization scope}), or arbitrarily select one rule when precedence is unclear (\emph{conditional precedence}).
Linguistic-layer types produce ``misparse-and-apply'' failures: agents adopt one reading of a scope ambiguity (\emph{scopal}) or misresolve a referential expression (\emph{coreferential}).
\emph{Lexical} ambiguity is a surprising exception---agents treat word-sense ambiguity as an information gap (assumption-based action), not a misparse.

\paragraph{Authorization Scope Directionality.}
Among authorization-scope violations, 69.2\% involve unauthorized escalation (granting excessive permissions) and 0\% involve over-restriction, a highly significant directional bias (binomial $p$\,=\,5.5$\times$10$^{-5}$) consistent across 4/5 models.

\paragraph{Inter-Coder Reliability.}
We report inter-coder $\kappa$\,=\,0.517 (moderate), with percentage agreement of 65.6\%.
The moderate $\kappa$ reflects fine-grained category boundary disagreements, particularly at the unauthorized-escalation/inappropriate-action boundary.
The three key findings above---Ambiguity Blindness rates, layer mechanism divergence, and authorization-scope directionality---are robust binary or majority-level metrics unaffected by the moderate fine-grained $\kappa$.

% ============================================================
% 7. DISCUSSION
% ============================================================
\section{Discussion}

\paragraph{Differential Nocuity.}
Berry and Kamsties' framework distinguishes nocuous from innocuous ambiguity in human requirements engineering.
Our results show that \emph{all} policy ambiguity types are nocuous for LLM agents, but the degree of nocuity varies by both type and model.
We term this \emph{Differential Nocuity}: the binary effect is robust ($+$34.0pp on $\tau^2$-bench), but effect size is moderated by benchmark strictness (OR\,=\,15.1 with LLM judge vs.\ OR\,=\,3.31 with rule-based verifier on SOPBench) and domain (healthcare $\Delta$\,=\,16.2pp $<$ banking $\Delta$\,=\,20.4pp $<$ customer service $\Delta$\,=\,34.0pp).

\paragraph{Systematic Nocuity.}
Our data reveals a phenomenon beyond Berry and Kamsties' binary: in 27\% of ambiguous episodes, all models converge on the same wrong interpretation.
This \emph{systematic nocuity} is more dangerous than divergent nocuity because it appears correct in inter-model agreement checks and is invisible until compared against ground truth.

\paragraph{Safety Training Selectivity.}
The model$\times$type interaction reveals that safety training is not uniformly effective across ambiguity types.
Claude's 12\% authorization-scope violation rate versus GPT-5.4's 64\% suggests that Anthropic's RLHF specifically addresses authorization-boundary reasoning.
Yet all models fail similarly on \emph{incompleteness} (46--62\%), indicating that safety training cannot compensate for literally absent information---a fundamental limitation.

\paragraph{Practical Implications.}
For policy authors: \emph{incompleteness} should be the highest-priority target for disambiguation, as it is universally dangerous and unmitigatable by model selection.
For model developers: authorization-scope handling is a differentiable safety dimension where targeted training demonstrably helps.
For deployment teams: ambiguity detection must be externalized---agents' $<$5\% awareness rate means they cannot self-detect ambiguous policies.

% ============================================================
% 8. CONCLUSION
% ============================================================
\section{Conclusion}

We present the first causal study of how policy ambiguity types differentially affect LLM agent safety.
Across 2{,}997 episodes and five models, ambiguous policy clauses increase safety violations by 34.0pp, with a three-tier type hierarchy and dramatic model$\times$type interactions.
The binary effect replicates on SOPBench ($\Delta$\,=\,18.3pp) with domain moderation.
Ambiguity Blindness ($<$5\% awareness) means agents cannot self-detect dangerous policies.
Both type-aware policy auditing and model-specific risk assessment are needed; neither alone suffices.

% ============================================================
% LIMITATIONS
% ============================================================
\section*{Limitations}

Our study has three main limitations.
First, the primary evaluation uses $\tau^2$-bench (retail and airline domains); while SOPBench replication confirms the binary effect in healthcare and banking, effect size and type severity ordering are domain-moderated, and generalization to other domains remains to be tested.
Second, clause pairs are synthetically constructed, introducing slightly higher parse depth than naturally occurring ambiguities ($p$\,=\,0.04 vs.\ DRIP-R).
The matched-pair design requires this control, and three of four linguistic properties show no distributional difference.
Third, the layer effect (specification vs.\ linguistic, $\Delta$\,=\,6.08pp, $p$\,=\,0.017) is fragile under within-judge subgroup analysis (GPT-4.1 judged: $p$\,=\,0.059; GPT-5.4 judged: $p$\,=\,0.157) and should be interpreted as suggestive rather than confirmed (see Appendix~\ref{app:c3}).

% ============================================================
% REFERENCES (placeholder)
% ============================================================
\bibliography{references}
\bibliographystyle{acl_natbib}

\begin{thebibliography}{20}

\bibitem[\protect\citename{Berry and Kamsties}2003]{berry2003ambiguity}
Daniel~M. Berry and Erik Kamsties.
\newblock 2003.
\newblock Ambiguity in requirements specification.
\newblock In \emph{Perspectives on Software Requirements}, pages 7--44. Springer.

\bibitem[\protect\citename{Chantree \bgroup et al.\egroup}2006]{chantree2006identifying}
Francis Chantree, Bashar Nuseibeh, Anne de~Roeck, and Alistair Willis.
\newblock 2006.
\newblock Identifying nocuous ambiguities in natural language requirements.
\newblock In \emph{Proceedings of the 14th IEEE International Requirements Engineering Conference (RE)}, pages 59--68.

\bibitem[\protect\citename{Ferrari \bgroup et al.\egroup}2019]{ferrari2019nlp}
Alessio Ferrari, Giorgio~O. Spagnolo, and Stefania Gnesi.
\newblock 2019.
\newblock {NLP} for requirements engineering: Tasks, techniques, tools, and technologies.
\newblock In \emph{Proceedings of the IEEE 27th International Requirements Engineering Conference (RE)}, pages 422--423.

\bibitem[\protect\citename{Li \bgroup et al.\egroup}2024]{li2024drip}
Zhehao Li, Xingjian Bi, and others.
\newblock 2024.
\newblock {DRIP-R}: Disentangling ambiguity types in requirements interpretation by {LLMs}.
\newblock In \emph{Proceedings of EMNLP 2024}.

\bibitem[\protect\citename{Li \bgroup et al.\egroup}2025]{li2025sopbench}
Zhehao Li and others.
\newblock 2025.
\newblock {SOPBench}: Evaluating {LLM} agents on standard operating procedures.
\newblock In \emph{Proceedings of ACL 2025}.

\bibitem[\protect\citename{Massey \bgroup et al.\egroup}2014]{massey2014evaluating}
Aaron~K. Massey, Richard~L. Rutledge, Annie~I. Ant\'on, and Peter~P. Swire.
\newblock 2014.
\newblock Identifying and classifying ambiguity for regulatory requirements.
\newblock In \emph{Proceedings of the 22nd IEEE International Requirements Engineering Conference (RE)}, pages 83--92.

\bibitem[\protect\citename{Sclar \bgroup et al.\egroup}2024]{sclar2024quantifying}
Melanie Sclar, Yejin Choi, Yulia Tsvetkov, and Alane Suhr.
\newblock 2024.
\newblock Quantifying language models' sensitivity to spurious features in prompt design.
\newblock In \emph{Proceedings of ICLR 2024}.

\bibitem[\protect\citename{$\tau^2$-bench}2025]{tau2bench2025}
Sierra.
\newblock 2025.
\newblock $\tau^2$-bench: A benchmark for multi-turn agent safety evaluation.

\end{thebibliography}

% ============================================================
% APPENDIX
% ============================================================
\appendix

\section{SOPBench Cross-Benchmark Replication}
\label{app:sopbench}

We replicate C1 on SOPBench \cite{li2025sopbench} using 100 clause pairs across healthcare (50) and banking (50) domains, covering four types (\emph{incompleteness}, \emph{authorization\_scope}, \emph{lexical}, \emph{conditional\_precedence}).
With five models, this yields 1{,}000 target episodes (982 valid).
SOPBench uses a rule-based verifier rather than an LLM judge.

Table~\ref{tab:sopbench_c1} summarizes the results.
The binary effect replicates ($\Delta$\,=\,18.3pp, OR\,=\,3.31, $p<$0.0001) but is substantially weaker than the $\tau^2$-bench effect (34.0pp).
Three factors contribute to the gap: (1)~the rule-based verifier is stricter than LLM-based judging, estimated to account for 6--8pp; (2)~domain moderation, with healthcare showing lower effects (16.2pp) than banking (20.4pp); and (3)~higher unambiguous baselines in SOPBench (11.2\% vs.\ 6.9\%), compressing the delta.

\begin{table}[t]
\centering
\small
\begin{tabular}{llcccc}
\toprule
\textbf{Benchmark} & \textbf{Domain} & \textbf{$n$} & \textbf{$\Delta$pp} & \textbf{OR} & \textbf{$p$} \\
\midrule
$\tau^2$-bench & retail+airline & 2997 & 34.0 & 15.1 & $<$.0001 \\
SOPBench & healthcare+bank & 982 & 18.3 & 3.31 & $<$.0001 \\
\quad healthcare & -- & 487 & 16.2 & 3.06 & $<$.001 \\
\quad bank & -- & 495 & 20.4 & 3.58 & $<$.0001 \\
\bottomrule
\end{tabular}
\caption{Cross-benchmark C1 replication. The binary effect replicates on SOPBench with reduced magnitude.}
\label{tab:sopbench_c1}
\end{table}

Per-model SOPBench results: GPT-5.4 $\Delta$\,=\,22.5pp, GPT-4.1 $\Delta$\,=\,19.8pp, Qwen3 $\Delta$\,=\,18.4pp, DeepSeek $\Delta$\,=\,18.7pp, Claude $\Delta$\,=\,12.1pp.
Claude remains the most robust model across both benchmarks.

Type severity ranking on SOPBench: \emph{incompleteness} (39.2\%) $>$ \emph{lexical} (33.8\%) $>$ \emph{conditional\_precedence} (23.5\%) $>$ \emph{authorization\_scope} (21.5\%).
Incompleteness remains the top type; authorization scope drops from rank 3 ($\tau^2$-bench) to rank 4 (SOPBench), reflecting the more explicit permission structures in healthcare and banking SOPs.
Spearman $\rho$\,=\,0.40 ($p$\,=\,0.60) between the two rankings---partial alignment but not statistically significant with only four testable types.

\section{C3: Layer Effect (Exploratory)}
\label{app:c3}

Specification-layer types show higher aggregate violation rates than linguistic-layer types (44.0\% vs.\ 37.9\%, $\Delta$\,=\,6.08pp, $p$\,=\,0.017, $V$\,=\,0.062).
However, within-judge subgroup analysis reveals fragility: for episodes judged by GPT-4.1, $\Delta$\,=\,8.0pp ($p$\,=\,0.059); for episodes judged by GPT-5.4, $\Delta$\,=\,4.8pp ($p$\,=\,0.157).
The aggregate significance may partly reflect between-judge variance.

Per-model layer gaps vary dramatically: GPT-5.4 +14.0pp, GPT-4.1 +12.0pp, Claude +3.3pp, DeepSeek +2.0pp, Qwen3 $-$0.9pp (reversed).
The aggregate layer effect is driven by GPT models' high authorization-scope rates, making the model$\times$type interaction (C4) more fundamental than the layer main effect.
We present this finding as suggestive evidence that specification-layer ambiguity tends to be more dangerous, while noting that the pattern is model-dependent and not robust under subgroup analysis.

\section{Ambiguity Classifier Proof of Concept}
\label{app:classifier}

We train a DeBERTa-v3-large binary ambiguity detector and compare it with GPT-4.1 few-shot classification (Table~\ref{tab:classifier}).
Binary detection is viable (F1 $\approx$ 0.80), but six-type classification is weak, especially for \emph{incompleteness} (F1\,=\,0.31--0.34): absence of information provides no positive signal for detection.
\emph{Conditional precedence} is most detectable (F1\,=\,0.91--0.94) due to strong structural conflict signals.

\begin{table}[t]
\centering
\small
\begin{tabular}{lccc}
\toprule
\textbf{Method} & \textbf{Bin. F1} & \textbf{6-type F1} & \textbf{Worst type} \\
\midrule
DeBERTa-v3 & 0.79 & 0.55 & incomp. (0.31) \\
GPT-4.1 few-shot & 0.806 & 0.632 & incomp. (0.34) \\
\bottomrule
\end{tabular}
\caption{Ambiguity classifier comparison. Binary detection viable; type classification limited by absence-based types.}
\label{tab:classifier}
\end{table}

\section{TOST Equivalence Analysis}
\label{app:tost}

\begin{table}[t]
\centering
\small
\begin{tabular}{lccc}
\toprule
\textbf{Margin} & \textbf{Equiv.} & \textbf{Non-equiv.} & \textbf{Interpretation} \\
\midrule
$\pm$15pp & 8/15 & 7/15 & types not interchangeable \\
$\pm$10pp & 4/15 & 11/15 & most pairs distinguishable \\
$\pm$5pp & 1/15 & 14/15 & nearly all distinguishable \\
\bottomrule
\end{tabular}
\caption{TOST pairwise equivalence testing for C2. Even at the generous $\pm$15pp margin, 7/15 type pairs are significantly non-equivalent.}
\label{tab:tost}
\end{table}

\end{document}
